% Standalone problem document, generated from the adjacent README.md.
% Compile with XeLaTeX. Mathematical content is maintained in Markdown.
\documentclass[11pt,a4paper]{article}
\usepackage[fontsize=10.5pt]{scrextend}
\usepackage[margin=22mm,headheight=15pt,headsep=9mm,footskip=11mm]{geometry}
\usepackage{fontspec}
\setmainfont{texgyrepagella}[Extension=.otf,UprightFont=*-regular,BoldFont=*-bold,ItalicFont=*-italic,BoldItalicFont=*-bolditalic]
\setsansfont{texgyreheros}[Extension=.otf,UprightFont=*-regular,BoldFont=*-bold,ItalicFont=*-italic,BoldItalicFont=*-bolditalic]
\setmonofont{texgyrecursor}[Extension=.otf,UprightFont=*-regular,BoldFont=*-bold,ItalicFont=*-italic,BoldItalicFont=*-bolditalic,Scale=MatchLowercase]
\usepackage{amsmath,amssymb}
\usepackage{unicode-math}
\setmathfont{texgyrepagella-math.otf}
\usepackage{xcolor,microtype,enumitem,needspace,fancyhdr}
\usepackage{bookmark}
\definecolor{ink}{HTML}{17344B}
\definecolor{accent}{HTML}{197D80}
\definecolor{muted}{HTML}{596773}
\hypersetup{colorlinks=true,linkcolor=accent,urlcolor=accent,pdftitle={TR-11 - Generic
identifiability of tensors at strictly subcritical
ranks},pdfauthor={Open Problems in Numerical Linear Algebra}}
\urlstyle{same}
\setlength{\parindent}{0pt}
\setlength{\parskip}{5pt plus 1pt minus 1pt}
\linespread{1.04}
\setlength{\emergencystretch}{3em}
\widowpenalty=10000
\clubpenalty=10000
\displaywidowpenalty=10000
\allowdisplaybreaks[1]
\setlist{nosep,leftmargin=1.5em}
\providecommand{\tightlist}{\setlength{\itemsep}{0pt}\setlength{\parskip}{0pt}}
\providecommand{\pandocbounded}[1]{#1}
\pagestyle{fancy}
\fancyhf{}
\fancyhead[L]{\sffamily\footnotesize\color{muted}OPEN PROBLEMS IN NUMERICAL LINEAR ALGEBRA}
\fancyhead[R]{\sffamily\bfseries\footnotesize\color{accent}TR-11}
\fancyfoot[L]{\parbox[b]{0.9\textwidth}{\sffamily\fontsize{8}{10}\selectfont\color{muted}Editorial ratings; a bounded search cannot certify that a problem remains unsolved.\\Literature check: 2026-09-08}}
\fancyfoot[R]{\sffamily\footnotesize\color{muted}\thepage}
\renewcommand{\headrulewidth}{0.3pt}
\renewcommand{\headrule}{\hbox to\headwidth{\color{accent}\leaders\hrule height \headrulewidth\hfill}}
\makeatletter
\renewcommand{\section}{\Needspace{5\baselineskip}\@startsection{section}{1}{0pt}{12pt plus 2pt minus 2pt}{4pt}{\sffamily\bfseries\large\color{ink}}}
\renewcommand{\subsection}{\Needspace{4\baselineskip}\@startsection{subsection}{2}{0pt}{9pt plus 1pt minus 1pt}{3pt}{\sffamily\bfseries\normalsize\color{ink}}}
\makeatother
\setcounter{secnumdepth}{0}
\begin{document}
{\sffamily\small\color{muted}Tensor computations\par}
\vspace{3pt}
{\raggedright\hyphenpenalty=10000\exhyphenpenalty=10000\sffamily\bfseries\fontsize{21}{24}\selectfont\color{ink}Generic
identifiability of tensors at strictly subcritical ranks\par}
\vspace{7pt}
{\sffamily\small\textbf{Difficulty:} extreme\quad\textbf{Importance:} interesting
to the community\par}
\vspace{4pt}
{\color{accent}\hrule height 0.8pt}
\vspace{6pt}
\subsection{Statement}\label{statement}

Let \(d\ge3\), \(n_1\ge\cdots\ge n_d\ge2\), and define \[
r_*={\prod_{i=1}^d n_i\over 1+\sum_{i=1}^d(n_i-1)},\qquad
b=\prod_{i=2}^d n_i-\sum_{i=2}^d(n_i-1).
\] For every integer \(1\le r<r_*\), a generic rank-\(r\) tensor in
\(\bigotimes_{i=1}^d\mathbb C^{n_i}\) is conjectured to have a unique
expression as a sum of \(r\) nonzero rank-one tensors, except in these
cases:

\begin{enumerate}
\def\labelenumi{\arabic{enumi}.}
\tightlist
\item
  \(n_1>b\) and \(r\ge b\).
\item
  Format \((4,4,3)\), rank \(5\).
\item
  Format \((n,n,2,2)\), rank \(2n-1\), with \(n\ge2\).
\item
  Format \((4,4,4)\), rank \(6\).
\item
  Format \((6,6,3)\), rank \(8\).
\item
  Format \((2,2,2,2,2)\), rank \(5\).
\end{enumerate}

A rank-one summand is \(v_1\otimes\cdots\otimes v_d\). Uniqueness
identifies decompositions whose tensor summands differ only by
permutation. Generic means outside a proper Zariski-closed subset of the
Zariski closure of tensors of rank at most \(r\), restricted to tensors
of actual rank \(r\). The number \(r_*\) is a rational dimension-count
value, not an assumed actual generic rank.

\subsection{Relevance}\label{relevance}

The conjecture specifies the range in which a generic low-rank tensor
model has uniquely recoverable factors, an essential precondition for
numerical CP decomposition and component interpretation.

\subsection{References}\label{references}

\begin{enumerate}
\def\labelenumi{\arabic{enumi}.}
\tightlist
\item
  L. Chiantini, G. Ottaviani, and N. Vannieuwenhoven, \emph{Effective
  criteria for specific identifiability of tensors and forms}.
  \href{https://arxiv.org/pdf/1609.00123}{Primary preprint}, §3,
  Conjecture 6 and Remark 7, p.7. The author journal manuscript numbers
  these Conjecture 3.4 and Remark 3.5.
\item
  The same authors, \emph{An Algorithm for Generic and Low-Rank Specific
  Identifiability of Complex Tensors}, SIAM J. Matrix Anal. Appl. 35
  (2014), 1265--1287.
  \href{https://epubs.siam.org/doi/10.1137/140961389}{DOI};
  \href{https://arxiv.org/abs/1403.4157}{primary preprint}, Theorem 1.1,
  for the verified bounded-format classification.
\item
  A. Massarenti and M. Mella, \emph{Bronowski's conjecture and the
  identifiability of projective varieties}.
  \href{https://arxiv.org/pdf/2210.13524v3}{January 2024 primary
  revision}, abstract and §§3--4, for later conditional links between
  identifiability and nondefectivity.
\end{enumerate}

\subsection{Status check --- 2026-09-08}\label{status-check-2026-09-08}

Reference 1 records verification for \(\prod n_i\le15000\). The later
Massarenti--Mella results reduce broad identifiability questions to
secant nondefectivity under hypotheses; they are not a full
classification of these formats. Searches
\textit{"generic identifiability" "conjecture" "2026" tensor},
\textit{"Chiantini" "Vannieuwenhoven" "conjecture" "2025"}, and
\textit{"Chiantini" "Vannieuwenhoven" "conjecture" "2026"} located no
full resolution. Status is bounded by those searches.
\end{document}
